\section{Proofs}\label{app:proofs}

This appendix gives the proofs that \cref{sec:requests} sketches. We use the core
calculus of \cref{sec:providers} and the semantics of \cref{sec:couplings}: given a tape
$\rho$, an execution is deterministic, and $\langle c, \sigma, h\rangle \Downarrow_\rho
\langle \sigma', h' \rangle$ runs $c$ from store $\sigma$ and history $h$. A call
$x \leftarrow \mathsf{call}\ r$ evaluates $r$ in $\sigma$ to a request $q$, computes the
key $k$ of the event from $h$ under the scheme, obtains $(\mathit{resp}, o) = \Pi(m, q,
\rho(k))$, binds $x$ to $\mathit{resp}$ and appends $\mathsf{call}(m, q, \mathit{resp}, o)$.
A retry runs its body, lets $V$ decide from the attempt's transcript and a draw keyed by
it, appends $\mathsf{attempt}(b)$, and stops if $b$ holds or no attempts remain.

\begin{proof}[Proof of \cref{lem:equal}]
By induction on the length of the common signature, and within a node on its nesting.
Throughout, the two executions start from the same history, and we show that they
extend it by the same events and bind equal values at corresponding positions.
\begin{itemize}[leftmargin=*]
\item \emph{Call.} The pieces are equal. A $\ptext$ piece denotes the same text in both
  executions; a $\pvar(b)$ piece denotes $\sigma_1(b) = \sigma_2(b)$ by hypothesis; a
  $\pres(p)$ piece denotes the response bound at position $p$, which is equal in both
  executions by the induction hypothesis. So both evaluate the request to the same text
  $q$ for the same model identity $m$. The histories so far are equal, so the scheme,
  whose keys are functions of the history and the request, computes the same key $k$;
  the tape gives the same draw; and $\Pi(m, q, \rho(k))$ gives the same response and
  count. Both histories are extended by the same event, and both bindings receive the
  same response.
\item \emph{Retry.} The bounds are equal. Each attempt runs a body with equal signatures
  from equal histories, so by the induction hypothesis it produces equal transcripts.
  The validator reads the same transcript and, the histories being equal, the same draw,
  so both executions decide alike and stop after the same attempt.
\item \emph{Public branch.} The guard terms are equal and denote equal values, a $\pvar$
  piece by hypothesis and a $\pres$ piece by induction, so both executions take the same
  arm, whose signatures are equal.
\item \emph{Resolved branch.} $\sigv{v_i}$ inlines the arm that $v_i$ selects, and
  $v_i$ is the outcome vector of $\sigma_i$, so the inlined arm is the one execution $i$
  takes. The inlined nodes are then compared like any others. That a store's own outcome
  vector selects the arm the store takes is \code{resolve\_exact} in the Coq
  development; it needs the side condition that no statement before the branch
  rebinds a secret, which the type system guarantees.\qedhere
\end{itemize}
\end{proof}

\begin{proof}[Proof of \cref{thm:leakage}]
Let $\mathit{cls}(\sigma_S)$ be the class of $\sigv{v(\sigma_S)}(c)$ among the at most $k(c)$
distinct signatures. Fix public inputs $\sigma_P$ and a tape $\rho$. By \cref{lem:equal},
two secret assignments with the same class produce the same history; signatures are
symbolic in the public inputs, so this holds for every $\sigma_P$. Outputs and effects
are functions of public data and responses: the \textsc{Output} and \textsc{Emit} rules
of \cref{fig:rules} reject any output or effect under a secret program counter, and any
secret content reaching them. Hence one invocation's observation is
$G(\mathit{cls}(\sigma_S), \sigma_P, \rho)$ for some function $G$.

Now let an adversary choose the public inputs of each invocation from the observations
of earlier invocations and its own coins $\alpha$, with fresh tapes $\rho_1, \rho_2,
\dots$ By induction on the number of invocations, the whole sequence of observations is
a function of $\mathit{cls}(\sigma_S)$, the tapes and $\alpha$: the $j$-th public input is
a function of the first $j-1$ observations and $\alpha$, and the $j$-th observation is
$G(\mathit{cls}(\sigma_S), \sigma_P^{(j)}, \rho_j)$. The tapes and $\alpha$ are independent
of the secrets. The channel from the secrets to the observations therefore factors as a
cascade whose first stage is the deterministic map $\mathit{cls}$ onto at most $k(c)$
values. The min-entropy leakage of a cascade is at most that of its first
stage~\cite{espinoza2012min}, and for every prior the min-entropy leakage of a channel
with at most $k$ possible outputs is at most $\log_2 k$~\cite[Theorem
1]{doychev2013cacheaudit}. Taking the maximum over priors gives min-capacity at most
$\log_2 k(c)$.
\end{proof}

\begin{proof}[Proof of \cref{thm:observers}]
Consider a maximal run of calls outside any retry body, none of which references
another in the run. Every request in the run is determined by pieces from outside the
run, so both executions send the same multiset of requests in the run (for the bill
normal form) or the same sequence for each provider (for the provider normal form).
Under the request scheme a call's draw is keyed by its model, its request and the
number of earlier calls with the same request, not by its position; a stable sort keeps
equal requests in their order of execution, so the $i$-th copy of a request has the same
occurrence index in both executions and receives the same response. Totals per model
are independent of order, and each provider's own sequence is preserved by the stable
sort. The $\pres$ references are rewritten to follow the moved calls, so later pieces
denote the same responses. Control nodes are compared exactly and sit between runs, so
the induction of \cref{lem:equal} continues past them. Calls inside a retry body are not
moved, because the validator reads the attempt's transcript in order.
\end{proof}

\begin{proof}[Proof of \cref{thm:completeness}]
Choose public inputs that are pairwise distinct strings over a character that occurs in
no constant of the program, and a provider whose responses are fresh strings that encode
the model, the request and the draw, with a validator that always rejects, so every
retry runs to its bound. Under this provider the rendering of a signature as a history
is injective: distinct sequences of pieces give distinct request texts, because the
public inputs and responses are distinct and fresh; distinct retry bounds give distinct
numbers of attempts; and a call present in one signature and absent in the other changes
the length of the history. The hypothesis on public guards provides inputs and responses
that drive both executions to the first difference, before which their histories agree by
\cref{lem:equal}; at the difference the histories differ, so the trace distinguishes the
two secrets. For $k(c)$ classes, the same provider makes the trace an injective function
of the class, a deterministic channel with $k(c)$ distinct outputs, whose min-entropy
leakage under the uniform prior on the classes is $\log_2
k(c)$~\cite[Theorem 1]{smith2009foundations}.
\end{proof}

\begin{proof}[Proof of \cref{prop:pc}]
Under the program-counter rule of Volpano et al.~\cite{volpano1996sound}, an event
observable at the public level may not occur under a secret program counter, and every
call inside a secret-guarded branch occurs under one; so every such branch is rejected.
If calls are unobservable, the rule places no constraint on them, and \cref{fig:triage},
whose only secret-dependent behaviour is a call, is accepted.
\end{proof}

\section{The real-workflow corpus}\label{app:corpus}

\Cref{tab:ports} lists the 30 ports. The adaptations are those of the protocol: A1, a
whole response is passed where the source extracts a field; A2, a multi-way choice
becomes a cascade of yes/no calls; A3, an early exit decided by a response is dropped;
A4, a list that the model filters is passed unfiltered; A6, a stand-in replaces a value
computed inside a branch as the single output. Each adaptation is recorded in the
corpus manifest with whether the certified bound still bounds the source. The 22
excluded candidates, each with its category and reason, are listed in the exclusion log
of the artifact (\file{bench/real/EXCLUSIONS.md}).

\begin{table}[p]
\centering
\caption{The 30 ported workflows (E11). Codes are those reported on the port without
annotations; endorsements are those of the accepted variant. The output bound and the
guaranteed input bound are in tokens; ``--'' means that a model's tokenizer is unpublished.
The last column gives AgentDojo's recorded successful attacks within the task's plan, out of
all successful attacks. Prefixes: \code{ac} Anthropic cookbook, \code{lg} LangGraph,
\code{ad} AgentDojo.}
\label{tab:ports}
\scriptsize\setlength{\tabcolsep}{3pt}
\begin{tabular}{@{}llllrlrrl@{}}
\toprule
Port & Split & Verdict & Codes & End. & Adapt. & Output & Input (guar.) & Attacks \\
\midrule
% Generated by figures/data.py from bench/results/real.json. Do not edit.
\code{ac-chain} & dev & accept & -- & 0 & -- & 16,384 & -- & -- \\
\code{ac-parallel} & dev & accept & -- & 0 & A6 & 16,384 & -- & -- \\
\code{ac-route} & held-out & accept & W237 & 0 & A1 A2 & 16,384 & -- & -- \\
\code{ac-evaluator-optimizer} & held-out & accept & -- & 0 & A1 A3 & 24,576 & -- & -- \\
\code{lg-agent-simulation} & held-out & accept & -- & 0 & A3 & 32,768 & 14,681,956 & -- \\
\code{lg-code-assistant} & dev & reject & E233 E235 & 8 & A1 A6 & 12,288 & -- & -- \\
\code{lg-extraction-retries} & held-out & accept & -- & 0 & A1 A3 & 12,288 & -- & -- \\
\code{lg-crag} & held-out & reject & E233 E235 & 4 & A4 & 73,728 & 642,244 & -- \\
\code{lg-reflection} & dev & accept & -- & 0 & -- & 229,376 & 17,262,535 & -- \\
\code{lg-self-discover} & held-out & accept & -- & 0 & -- & 16,384 & 1,611,298 & -- \\
\code{ad-banking-0} & held-out & reject & E233 & 1 & -- & 8,192 & 17,474 & 0/7 \\
\code{ad-banking-1} & held-out & accept & -- & 0 & -- & 4,096 & 8,716 & 0/7 \\
\code{ad-banking-2} & held-out & reject & E233 & 1 & -- & 8,192 & 33,908 & 0/6 \\
\code{ad-banking-3} & held-out & reject & E233 & 1 & -- & 8,192 & 17,802 & 0/7 \\
\code{ad-banking-4} & held-out & reject & E233 & 1 & -- & 8,192 & 17,478 & 0/7 \\
\code{ad-slack-0} & held-out & accept & -- & 0 & -- & 4,096 & 8,726 & 0/5 \\
\code{ad-slack-1} & dev & reject & E233 & 1 & -- & 12,288 & 42,665 & 0/5 \\
\code{ad-slack-2} & held-out & reject & E233 & 1 & -- & 8,192 & 17,518 & 0/5 \\
\code{ad-slack-3} & held-out & reject & E233 & 1 & -- & 8,192 & 17,516 & 0/5 \\
\code{ad-slack-4} & held-out & reject & E233 & 1 & -- & 12,288 & 42,977 & 0/5 \\
\code{ad-travel-0} & held-out & reject & E235 & 1 & -- & 8,192 & 17,926 & 0/3 \\
\code{ad-travel-1} & held-out & reject & E233 E235 & 1 & -- & 12,288 & 51,861 & 0/2 \\
\code{ad-travel-2} & dev & reject & E233 & 2 & -- & 24,576 & 111,748 & 0/1 \\
\code{ad-travel-3} & dev & reject & E233 & 2 & -- & 16,384 & 85,584 & 0/2 \\
\code{ad-travel-4} & held-out & reject & E233 & 2 & -- & 20,480 & 94,919 & 0/0 \\
\code{ad-workspace-0} & held-out & accept & -- & 0 & -- & 4,096 & 8,800 & 0/6 \\
\code{ad-workspace-1} & held-out & accept & -- & 0 & -- & 4,096 & 8,793 & 0/4 \\
\code{ad-workspace-2} & held-out & accept & -- & 0 & -- & 4,096 & 16,943 & 0/3 \\
\code{ad-workspace-3} & dev & accept & -- & 0 & -- & 4,096 & 8,746 & 0/3 \\
\code{ad-workspace-5} & held-out & accept & -- & 0 & -- & 4,096 & 8,817 & 0/5 \\
\bottomrule

\end{tabular}
\end{table}
